\documentclass[11pt]{article}
\usepackage[margin=1in]{geometry}
\usepackage{hyperref}
\usepackage{listings}
\title{EV3 Line Follower -- Assignment 2 Guide}
\date{30 September 2026}
\begin{document}\maketitle
\section{Background}
This work controls a LEGO EV3 robot on yellow tape over a dark surface. A reflected-light color sensor supplies three states (WHITE, MIDDLE, BLACK); a Q table selects forward, left, or right actions. An infrared sensor detects obstacles. Q learning handles local steering; a junction route needs an explicit policy or extra sensing.
\section{Files and duties}
\begin{description}
\item[\texttt{safe\_main.py}] Final controller: motors A/D, color S1, IR S4, Q-table line following, Center stop, obstacle beep/reverse/turn/path reacquisition, and training mode.
\item[\texttt{guided\_calibration.py}] Records 20 reflection samples for dark background, yellow tape, and tape edge; writes \texttt{calibration.json}.
\item[\texttt{check\_motor\_direction.py}] Pulses A and D with wheels lifted to verify wiring and direction.
\item[\texttt{make\_qtable\_ev3.py}] Converts the readable 18-entry seed to \texttt{q\_table.pkl}.
\item[\texttt{q\_table\_simulated.txt}] Human-readable seed table.
\item[\texttt{EV3\_Line\_Follower\_Complete\_Guide.pdf}] Reference guide; boot files are retained but never flashed.
\end{description}
\section{Code and training}
The controller class reads reflection, maps it to WHITE/MIDDLE/BLACK, executes motor commands and stops on Center or an obstacle, and brakes both motors during cleanup. \texttt{best\_action} selects the highest Q value for $(mode,state)$. On an obstacle the robot beeps, reverses, pivots without RL, accepts stable WHITE or MIDDLE as a recovered path, beeps again, and resumes. Training observes a transition, gives +10 for MIDDLE and -10 otherwise, applies $Q\leftarrow Q+\alpha(r+\gamma Q'_{max}-Q)$, and saves \texttt{q\_table\_candidate.pkl}. Keep a copy of the seed table before training.
\section{Complete commands}
Use the current Bluetooth address shown by the EV3. Open the terminal in the folder containing these files. The commands below use WSL syntax.
\begin{lstlisting}
EV3_IP=169.254.73.240
# Computer (WSL)

ping $EV3_IP
ssh robot@$EV3_IP                    # password: maker

# EV3 SSH shell
mkdir -p /home/robot/assignment2
exit

# Computer: upload
scp safe_main.py guided_calibration.py check_motor_direction.py make_qtable_ev3.py q_table_simulated.txt robot@$EV3_IP:/home/robot/assignment2/
ssh robot@$EV3_IP
cd /home/robot/assignment2

# EV3: build, calibrate, motor test
brickrun -r -- pybricks-micropython make_qtable_ev3.py
brickrun -r -- pybricks-micropython guided_calibration.py
brickrun -r -- pybricks-micropython check_motor_direction.py

# Normal run (edit safe_main.py so TRAINING = False)
brickrun -r -- pybricks-micropython safe_main.py
# Center starts and stops

# Training (backup, edit TRAINING = True, then run)
cp q_table.pkl q_table_seed_backup.pkl
nano safe_main.py
brickrun -r -- pybricks-micropython safe_main.py
# candidate: q_table_candidate.pkl; Center starts/stops

# Retrieve after exit
exit
scp robot@$EV3_IP:/home/robot/assignment2/calibration.json .
scp robot@$EV3_IP:/home/robot/assignment2/q_table_candidate.pkl .
scp robot@$EV3_IP:/home/robot/assignment2/*.log .
\end{lstlisting}
\section{Motor EIO error}
For \texttt{OSError: [Errno 5] EIO} on line 9, inspect ports:
\begin{lstlisting}
for d in /sys/class/tacho-motor/*; do echo "$d"; cat "$d/address"; cat "$d/driver_name"; done
\end{lstlisting}
Both \texttt{outA} and \texttt{outD} must report \texttt{lego-ev3-l-motor}. Power off, reconnect cables, reboot, and test independently:
\begin{lstlisting}
brickrun -r -- pybricks-micropython -c "from pybricks.ev3devices import Motor; from pybricks.parameters import Port; Motor(Port.A); print('A OK')"
brickrun -r -- pybricks-micropython -c "from pybricks.ev3devices import Motor; from pybricks.parameters import Port; Motor(Port.D); print('D OK')"
\end{lstlisting}
\end{document}








